% =====================================================================================
%  Appendix A -- the proof of Theorem~\ref{thm:leakage} and of the closed form
%  Eq.~\eqref{eq:lambda}.  Items (1)-(5) are those of the Supplemental Material section
%  sm:leakage, moved here verbatim so that the main text carries the proof of its
%  central theorem; \label{eq:split} lives here, and the SM reaches it through xr-hyper.
%  main.tex inputs this file after sec_10.tex (which ends with the acknowledgments) and
%  before bibliography.tex, i.e. before sm_begin.tex, so pra_split.py places it in
%  main_pra.pdf.  \appendix itself is issued in main.tex, next to the \input.
% =====================================================================================
\section{Proof of the leakage certificate and of its closed form}
\label{app:proof}

This appendix proves Theorem~\ref{thm:leakage} and the second equality of Eq.~\eqref{eq:lambda}.

\begin{proof}
Write $u=R(z)\phi$ and $u_P=R_P(z)\phi_P$, extended by zero outside $\operatorname{ran}P$.

\emph{(1) The split.} Since $u_P\in\operatorname{ran}P$ and $(z-H_P)u_P=\phi_P$ there,
$(z-H)u_P=zu_P-PHu_P-QHu_P=\phi_P-QHu_P$; subtracting this from $(z-H)u=\phi$ gives
\begin{equation}
u-u_P\;=\;R(z)\bigl[\,\phi_Q\;+\;QH\,R_P(z)\phi_P\,\bigr].
\label{eq:split}
\end{equation}
This is the Feshbach--L\"owdin partition of the resolvent~\cite{feshbach1958,lowdin1962}, written as an
exact identity rather than as an expansion: the first source term is the part of the probe the subspace
never held, the second is the part of the Hamiltonian that carries amplitude out of it. Only $P^{2}=P$,
$P^{\dagger}=P$, $H=H^{\dagger}$ and $\eta>0$ are used; no property of the determinant basis, of the
ranking, or of the sampler enters anywhere below.

\emph{(2) Pointwise.}
$A-A_P=(\eta/\pi)(\lVert u\rVert+\lVert u_P\rVert)(\lVert u\rVert-\lVert u_P\rVert)$ and
$\bigl|\lVert u\rVert-\lVert u_P\rVert\bigr|\le\lVert u-u_P\rVert$, so
$|A-A_P|\le(\eta/\pi)(\lVert u\rVert+\lVert u_P\rVert)\,\lVert u-u_P\rVert$ at every $\omega$.

\emph{(3) Integration.} By Cauchy--Schwarz in $\omega$,
\begin{widetext}
\begin{equation*}
\lVert A-A_P\rVert_1\;\le\;\frac{\eta}{\pi}
\Bigl[\int\bigl(\lVert u\rVert+\lVert u_P\rVert\bigr)^{2}d\omega\Bigr]^{1/2}
\Bigl[\int\lVert u-u_P\rVert^{2}d\omega\Bigr]^{1/2}.
\end{equation*}
\end{widetext}
The first bracket is at most $(\pi/\eta)^{1/2}(\lVert\phi\rVert+\lVert\phi_P\rVert)$, by Minkowski in
$L^{2}(d\omega)$ and the two sum rules. For the second, apply Minkowski to Eq.~\eqref{eq:split}: its
first term obeys the sum rule with seed $\phi_Q$, contributing exactly
$(\pi/\eta)^{1/2}\lVert\phi_Q\rVert$; its second obeys
$\lVert R(z)v\rVert\le\eta^{-1}\lVert v\rVert$, since $H$ is self-adjoint and $\operatorname{Im}z=\eta$,
and the definition~\eqref{eq:lambda} then contributes at most $(\pi/\eta)^{1/2}\Lambda_P(\eta)/\eta$.
Multiplying the three factors gives the leakage branch of~\eqref{eq:twobudget}.

\emph{(4) The other two branches.} $A$ and $A_P$ are nonnegative with known integrals, so
$\lVert A-A_P\rVert_1\le\int A+\int A_P=\lVert\phi\rVert^{2}(1+w_P)$ under no hypothesis at all; and
$\bigl|\int A-\int A_P\bigr|\le\lVert A-A_P\rVert_1$ gives the lower bound
$\lVert\phi\rVert^{2}(1-w_P)$.

\emph{(5) The closed form.} $R_P(z)\phi_P=\sum_m c_m(z-\tilde E_m)^{-1}\ket{\tilde m}$, so
\begin{equation*}
\bigl\lVert QHR_P(z)\phi_P\bigr\rVert^{2}\;=\;\sum_{m,m'}
\frac{c_m\bar c_{m'}\,\braket{g_{m'},g_m}}{(z-\tilde E_m)(\bar z-\tilde E_{m'})}\,.
\end{equation*}
With $x=\omega+E_0$ the remaining integral is
$\int dx\,[(x-\tilde E_m+i\eta)(x-\tilde E_{m'}-i\eta)]^{-1}$, which has one pole in each half
plane; closing above picks up $2\pi i/[(\tilde E_{m'}-\tilde E_m)+2i\eta]$, and multiplying by
$\eta/\pi$ gives Eq.~\eqref{eq:lambda}.
\end{proof}
